\section{Fock空间与产生湮灭算符}\label{sec:fock}
上一节中我们完成了自由多粒子Hilbert空间的构造，
并将态按粒子数分为不同空间。
但仍有两个问题悬而未决：

\begin{enumerate}
\item \textbf{技术问题}：始终手动保持态的全对称或全反对称，
      在实际计算中极不方便。
      能否发展一种\emph{自动}保持对称性的代数工具？
\item \textbf{基本问题}：
      如果多粒子Hilbert空间上只有动量算符$\boldsymbol{p}$，
      我们只能构造对角的可观测量。
      特别地，有没有算符能在\emph{不同粒子数}的态上具有非零矩阵元，
      从而把不同空间联系起来？
\end{enumerate}

\noindent 引入\textbf{产生湮灭算符}将同时解决这两个问题。


为了描述粒子数可变的系统
（相对论中$E=mc^2$允许粒子对的产生与湮灭），
将所有粒子数空间的Hilbert空间做\textbf{直和}，
定义\textbf{Fock空间}：
\begin{equation}\label{eq:fockspace}
\boxed{\mathbb{F} = \bigoplus_{N=0}^{\infty}\mathbb{H}_N}
\end{equation}
其中$\mathbb{H}_0=\mathbb{C}$是真空态$|0\rangle$所在的一维空间，
$\mathscr{H}_N$是$N$粒子态的（对称化或反对称化的）张量积空间。

Fock空间的物理图像：
它是所有可能粒子数的Hilbert空间的直和——
$\mathbb{H}_0$描述真空，
$\mathbb{H}_1$描述单粒子态，
$\mathbb{H}_2$描述两粒子态，以此类推。
一个一般的Fock空间态可以是不同粒子数态的\textbf{叠加}——
这正是粒子可以产生和湮灭的数学基础。
\topictitle{产生算符}
在Fock空间上定义\textbf{产生算符}$a^\dagger(\boldsymbol{p},\sigma,n)$，
其作用是在系统中添加一个动量$\boldsymbol{p}$、自旋$\sigma$、种类$n$的粒子：
\begin{equation}\label{eq:creationdef}
a^\dagger(\boldsymbol{p},\sigma,n)\,
|\boldsymbol{p}_1,\sigma_1,n_1;\;\cdots;\;\boldsymbol{p}_N,\sigma_N,n_N\rangle
= |\boldsymbol{p},\sigma,n;\;\boldsymbol{p}_1,\sigma_1,n_1;\;\cdots;\;
\boldsymbol{p}_N,\sigma_N,n_N\rangle\,.
\end{equation}
从真空出发，反复施加产生算符即可构造任意多粒子态：
\begin{equation}\label{eq:fockbuild}
|\boldsymbol{p}_1,\sigma_1,n_1;\;\cdots;\;\boldsymbol{p}_N,\sigma_N,n_N\rangle
= a^\dagger(\boldsymbol{p}_1,\sigma_1,n_1)\cdots
  a^\dagger(\boldsymbol{p}_N,\sigma_N,n_N)\,|0\rangle\,.
\end{equation}

产生算符的引入在物理上是\textbf{必要的}：
在相对论能标下粒子可以被产生，
因此描述相互作用的算符\emph{必须}能够改变粒子数。
产生算符正是实现这一功能的最自然的数学工具。
\medskip


湮灭算符产生可以从多粒子态的正交关系\eqref{eq:multinorm}直接推导出来。以$a(\boldsymbol{p},\sigma,n)$左乘一个$N$粒子态，
利用\eqref{eq:multinorm}计算其与任意$(N{-}1)$粒子态的内积，可以得到：
\begin{equation}\label{eq:annihilate}
a(\boldsymbol{p},\sigma,n)\,|\boldsymbol{p}_1,\sigma_1,n_1;\;\cdots;\;
\boldsymbol{p}_N,\sigma_N,n_N\rangle
= \sum_{k=1}^{N}(\pm 1)^{k-1}\,
\langle\boldsymbol{p},\sigma,n|\boldsymbol{p}_k,\sigma_k,n_k\rangle\;
|\boldsymbol{p}_1,\ldots,\widehat{\boldsymbol{p}_k},\ldots,\boldsymbol{p}_N\rangle\,.
\end{equation}
其中$\widehat{\boldsymbol{p}_k}$表示省略第$k$个粒子，
正号取玻色子、负号取费米子。
物理含义很清楚：
$a(\boldsymbol{p})$在态中"寻找"动量为$\boldsymbol{p}$的粒子并将其\textbf{湮灭}；
如果态中不含这样的粒子，作用结果为零。
特别地：
\begin{equation}\label{eq:avac}
a(\boldsymbol{p},\sigma,n)\,|0\rangle = 0\,.
\end{equation}
即真空中没有任何粒子可以被湮灭。


\medskip

在一个态上先产生再湮灭，与先湮灭再产生，
结果一般不同——差别恰好编码了粒子的统计性质。
直接计算可以验证，对任意态都有：
\begin{tcolorbox}[colback=white!5,colframe=black!45]
\textbf{玻色子}和\textbf{费米子}的产生湮灭算符满足\textbf{对易关系}：
\begin{subequations}
      \begin{align}
      \big[a(\boldsymbol{p},\sigma,n),\;a^\dagger(\boldsymbol{p}',\sigma',n')\big]_\mp
      &= \delta^3(\boldsymbol{p}-\boldsymbol{p}')\,\delta_{\sigma\sigma'}\,\delta_{nn'}\,,
      \label{eq:bCR1}\\[4pt]
      \big[a(\boldsymbol{p},\sigma,n),\;a(\boldsymbol{p}',\sigma',n')\big]_\mp
      &= 0\,,\label{eq:bCR2}\\[4pt]
      \big[a^\dagger(\boldsymbol{p},\sigma,n),\;a^\dagger(\boldsymbol{p}',\sigma',n')\big]_\mp
      &= 0
\end{align}
\end{subequations}\label{eq:bCR3}
\end{tcolorbox}

\noindent 费米子的反对易关系\eqref{eq:bCR3}立即给出\textbf{Pauli不相容原理}：
\begin{equation}
a^\dagger(\boldsymbol{p},\sigma,n)\,a^\dagger(\boldsymbol{p},\sigma,n)\,|0\rangle
= 0\,,
\end{equation}
即不可能有两个全同费米子占据同一量子态。

\medskip

至此，本节开头的两个问题都已解决：

\begin{enumerate}
\item \textbf{对称化是自动的}：
      只需使用对易关系
      或反对易关系\eqref{eq:bCR1}
      玻色子的全对称态和费米子的全反对称态在代数上是自动保证的。

\item \textbf{不同粒子数空间被联系起来}：
      $a^\dagger$将$\mathscr{H}_N$映射到$\mathscr{H}_{N+1}$，
      $a$将$\mathscr{H}_N$映射到$\mathscr{H}_{N-1}$。
      因此，用$a$和$a^\dagger$的多项式可以构造
      Fock空间上\textbf{任意}算符的矩阵元——
      这将在第\ref{chap:scattering}章讨论相互作用时变得关键。
      \end{enumerate}